% !TEX root = projet_session.tex
% Document autonome : aucune image ni aucun préambule externe requis.
% Consignes du projet. Les nouvelles hypothétiques restent hors de ce fichier.
% Pondération du cours : 30 %, comme le projet précédent (Slides/S1/s1.tex).
\documentclass[11pt,letterpaper]{article}
\usepackage[margin=2.5cm,headheight=14pt]{geometry}
\usepackage[T1]{fontenc}
\usepackage[utf8]{inputenc}
\usepackage[french]{babel}
\usepackage[sfdefault,light]{FiraSans}
\usepackage{xcolor}
\usepackage{microtype}
\usepackage{booktabs,tabularx,array}
\usepackage{enumitem}
\usepackage{tcolorbox}
\usepackage{titlesec}
\usepackage{fancyhdr,lastpage}
\usepackage{hyperref}

% Identité visuelle reprise des documents d'exercices du cours.
\definecolor{HECnavy}{HTML}{002855}
\definecolor{HECgreen}{HTML}{26D07C}
\definecolor{HECcoral}{HTML}{FF585D}
\definecolor{HECtext}{HTML}{2D3748}
\definecolor{HECgray}{HTML}{64748B}
\definecolor{HEClightgray}{HTML}{F1F5F9}
\color{HECtext}
\hypersetup{colorlinks=true,urlcolor=HECnavy,linkcolor=HECnavy,
  pdftitle={Taux sous pression -- Prévoir la décision de la Banque du Canada et ses implications pour votre industrie},
  pdfauthor={Jean-Félix Brouillette},pdfsubject={ECON 50803 -- Automne 2026}}
\urlstyle{same}
\setlength{\parindent}{0pt}
\setlength{\parskip}{4pt}
\setlength{\emergencystretch}{1.5em}
\setlist[itemize]{leftmargin=1.4em,itemsep=3pt,topsep=3pt,parsep=0pt}
\setlist[enumerate]{leftmargin=1.6em,itemsep=4pt,topsep=3pt,parsep=0pt}
\newcolumntype{Y}{>{\raggedright\arraybackslash}X}
\newcommand{\aconfirmer}{\textcolor{HECcoral!80!black}{À confirmer}}

\titleformat{\section}{\Large\bfseries\color{HECnavy}}{\thesection.}{0.5em}{}
  [\vspace{1pt}{\color{HECgreen}\rule{2.2cm}{1.5pt}}]
\titlespacing*{\section}{0pt}{0pt}{9pt}
\titleformat{\subsection}{\normalsize\bfseries\color{HECnavy}}{}{0pt}{}
\titlespacing*{\subsection}{0pt}{7pt}{3pt}
\newtcolorbox{consigne}[1]{colback=HEClightgray,colframe=HECnavy,
  boxrule=0pt,leftrule=3pt,arc=4pt,left=9pt,right=9pt,top=6pt,bottom=6pt,
  title={#1},coltitle=HECnavy,fonttitle=\bfseries,
  attach title to upper={\par\smallskip},before skip=7pt,after skip=7pt}
\newtcolorbox{pointcle}[1]{colback=HECgreen!6,colframe=HECgreen,
  boxrule=0pt,leftrule=3pt,arc=4pt,left=9pt,right=9pt,top=6pt,bottom=6pt,
  title={#1},coltitle=HECnavy,fonttitle=\bfseries,
  attach title to upper={\par\smallskip},before skip=7pt,after skip=7pt}

\pagestyle{fancy}
\fancyhf{}
\renewcommand{\headrulewidth}{0pt}
\fancyhead[L]{\small\textcolor{HECgray}{Taux sous pression}}
\fancyhead[R]{\small\textcolor{HECgray}{Automne 2026}}
\renewcommand{\footrulewidth}{0.4pt}
\fancyfoot[L]{\footnotesize\textcolor{HECgray}{ECON 50803\enspace\textbar\enspace Environnement macroéconomique}}
\fancyfoot[R]{\footnotesize\textcolor{HECgray}{\thepage\,/\,\pageref{LastPage}}}
\fancypagestyle{premiere}{\fancyhead{}}

\begin{document}
\thispagestyle{premiere}
\begin{flushleft}
{\small\bfseries\color{HECnavy} HEC MONTRÉAL\enspace\textbar\enspace MBA}\par
\vspace{5pt}
{\fontsize{24}{28}\selectfont\bfseries\color{HECnavy} Taux sous pression}\par
\vspace{5pt}
{\large\color{HECnavy} Prévoir la décision de la Banque du Canada et ses implications pour votre industrie}\par
\vspace{5pt}
{\small\color{HECnavy} Projet de session --- Automne 2026}\par
{\small Jean-Félix Brouillette\enspace\textbar\enspace ECON 50803}\par
{\color{HECgreen}\rule{3cm}{2pt}}
\end{flushleft}
\vspace{-5pt}
{\small\textcolor{HECgray}{Version de travail\enspace\textbar\enspace Pondération : 30\,\% de la note finale du cours}}

\begin{consigne}{Votre mandat}
Avec votre \textbf{équipe attitrée du programme de MBA}, vous devez prévoir la décision de la Banque du Canada (BdC) du \textbf{9 décembre 2026} et expliquer ses implications pour une industrie dans laquelle travaille actuellement au moins un membre de votre équipe. Présentez votre analyse le \textbf{2 décembre 2026}, défendez-la lors d'une conférence de presse et réévaluez-la à la lumière d'une nouvelle économique hypothétique. Vous jouerez aussi le rôle de journalistes auprès d'une autre équipe.
\end{consigne}

\vspace{0.15cm}

\subsection{Pourquoi ce projet est-il utile à un gestionnaire?}
Les décisions d'affaires reposent sur des hypothèses concernant les conditions macroéconomiques. Un détaillant doit anticiper le pouvoir d'achat de ses clients; une entreprise de construction, les conditions de crédit; un exportateur, la demande étrangère et le taux de change. Ce projet vous entraînera à \textbf{expliciter ces hypothèses et à en évaluer les risques} avant de prendre une décision.

\subsection{Objectifs d'apprentissage}
À l'issue du projet, vous serez en mesure de :
\begin{enumerate}
  \item \textbf{Interpréter l'actualité économique avec esprit critique.} Vous approprier le vocabulaire macroéconomique pour lire et discuter de l'actualité avec aisance. Utiliser les indicateurs économiques pour situer une nouvelle parmi les grandes tendances et distinguer un signal pertinent du bruit économique ambiant.
  \item \textbf{Traduire l'analyse macroéconomique en décisions d'affaires.} Apprendre à utiliser le \textit{Rapport sur la politique monétaire} (RPM) de la Banque du Canada pour interpréter la conjoncture, anticiper les décisions de politique monétaire et dégager des risques et des occasions propres à votre industrie afin d'éclairer des décisions de gestion.
  \item \textbf{Défendre et adapter un jugement de gestion.} Présenter une argumentation claire, étayée et convaincante, expliciter vos hypothèses et répondre aux objections. Examiner le raisonnement d'une autre équipe et formuler des questions pertinentes. Réviser votre conclusion lorsque de nouvelles informations le justifient. Vos compétences de présentation et de persuasion seront évaluées.
\end{enumerate}

\clearpage

\begingroup
\setlength{\parskip}{2pt}
\setlist[enumerate]{leftmargin=1.6em,itemsep=1pt,topsep=1pt,parsep=0pt}
\section{Prévoir la décision de la Banque du Canada}
\textbf{Prenez comme point de départ le \href{https://www.banqueducanada.ca/publication/rpm/}{Rapport sur la politique monétaire du 28 octobre 2026}.} Repérez le diagnostic, les prévisions et les principaux risques présentés par la BdC, puis confrontez-les aux données publiées depuis, jusqu'à la date limite commune de remise des diapositives. Expliquez ce qui conforte ou modifie cette analyse en vue de la décision de décembre.

Votre présentation devra répondre aux questions suivantes :
\begin{enumerate}
  \item \textbf{Quelle décision prévoyez-vous?} Annoncez une variation précise du taux directeur : maintien, baisse ou hausse, avec son ampleur en points de base. Une baisse de 25 points de base correspond à une baisse de 0.25 point de pourcentage. Précisez votre degré de confiance (sans nécessairement attribuer une probabilité précise à votre prévision).
  \item \textbf{Quels éléments justifient votre prévision?} Appuyez-vous sur différents indicateurs macroéconomiques ou financiers pertinents. Expliquez ce qu'ils révèlent et pourquoi certains pèsent davantage dans votre diagnostic.
\end{enumerate}
La question porte sur ce que vous prévoyez que la BdC \textbf{fera} et non pas ce qu'elle \textbf{devrait faire}.

\vspace{0.2cm}

\section{Analyser les implications pour votre industrie}
\textbf{Choisissez une industrie dans laquelle travaille actuellement au moins un membre de l'équipe.} Appuyez-vous sur le \textit{Rapport sur la politique monétaire du 28 octobre 2026}, complété par des données sectorielles ou des exemples professionnels concrets, sans divulguer d'information confidentielle. Le RPM n'a pas à traiter directement de votre industrie : tirez-en les implications en explicitant vos hypothèses. Répondez aux quatre questions suivantes :
\begin{enumerate}
  \item \textbf{Quelles perspectives le contexte macroéconomique laisse-t-il entrevoir pour votre industrie?} Appuyez-vous sur \textbf{au moins deux éléments du RPM}, avec les pages ou figures correspondantes, et sur les données plus récentes pertinentes. Expliquez leurs implications pour votre industrie au cours des \textbf{12 mois suivant la présentation} : demande, coûts, emploi ou investissement, selon les éléments les plus pertinents.
  \item \textbf{Quelle particularité rend votre industrie plus ou moins sensible à la politique monétaire?} Identifiez une caractéristique concrète et expliquez comment elle renforce ou atténue les effets d'une variation des taux. Par exemple : des achats financés à crédit et faciles à reporter, des stocks importants financés à taux variable, ou une demande pour des produits essentiels difficiles à remplacer ou à reporter. Choisissez une caractéristique pertinente pour votre industrie et justifiez votre choix.
  \item \textbf{Comment la décision de taux que vous prévoyez influencerait-elle ces perspectives?} Expliquez par quel canal cette décision\footnote{Si vous prévoyez un maintien, expliquez ce que la poursuite des conditions de taux actuelles implique pour votre industrie. Les effets des décisions passées peuvent continuer à se transmettre.} toucherait votre industrie et dans quels délais. Distinguez son effet supplémentaire de celui des conditions financières déjà en place. Un effet limité ou nul est une conclusion valable s'il est justifié.
  \item \textbf{Quelle décision de gestion serait touchée?} Identifiez une décision précise et expliquez comment les perspectives décrites dans votre analyse pourraient l'influencer : calendrier d'un investissement, niveau des stocks, plan d'embauche ou choix de financement, par exemple.
\end{enumerate}
\par
\endgroup

\newpage
\section{Présenter, questionner et réagir}
Chaque équipe présente son analyse une fois et joue le rôle de journalistes une fois. Chaque passage dure \textbf{19 minutes, transition comprise}. Les durées incluent les questions et les réponses et doivent être respectées.

\begin{tabularx}{\textwidth}{@{}>{\raggedright\arraybackslash}p{3.1cm}p{1.4cm}Y@{}}
\toprule
\textbf{Étape} & \textbf{Durée} & \textbf{Travail attendu} \\
\midrule
Présentation & 7 min & Défendre la prévision, les éléments de preuve et les implications pour l'industrie. \\
\addlinespace[4pt]
Conférence de presse & 4 min & Les journalistes posent exactement deux questions sur la prévision de décision de la BdC; l'équipe présentatrice y répond. \\
\addlinespace[4pt]
Nouvelle et délibération & 3 min & Le professeur révèle la nouvelle. L'équipe présentatrice délibère et prépare sa réponse. \\
\addlinespace[4pt]
Réponse à la nouvelle & 4 min & Expliquer le verdict révisé ou maintenu, le degré de confiance et les conséquences pour l'industrie. \\
\addlinespace[4pt]
Transition & 1 min & Installer l'équipe suivante; les diapositives sont chargées avant la séance. \\
\bottomrule
\end{tabularx}

\textbf{Déroulement des trois heures :} ouverture (4 min), quatre passages (76 min), pause (20 min), quatre passages (76 min), synthèse (4 min). Les huit passages totalisent 152 minutes, transitions comprises.

\subsection{Votre rôle de journalistes}
Le \textbf{30 novembre}, le professeur communiquera l'ordre des présentations, l'équipe que vous questionnerez et ses diapositives remises. Les affectations seront tirées au sort selon un cycle : chaque équipe en questionne une autre et reçoit les questions d'une troisième.

Pour vous préparer, \textbf{visionnez des extraits des périodes de questions d'au moins deux conférences de presse de politique monétaire de la BdC} : \href{https://www.banqueducanada.ca/multimedia/conference-de-presse-rapport-sur-la-politique-monetaire-janvier-2026/}{28 janvier 2026}, \href{https://www.banqueducanada.ca/multimedia/conference-de-presse-rapport-sur-la-politique-monetaire-juillet-2025/}{30 juillet 2025} ou \href{https://www.banqueducanada.ca/medias/discours/diffusions/}{autres webdiffusions}. Observez comment les journalistes questionnent une donnée, mettent une hypothèse à l'épreuve et font préciser un raisonnement.

Vous devez poser \textbf{exactement deux questions sur la prévision de décision de la BdC et sa justification}, avant la nouvelle. Choisissez librement vos angles : contradiction apparente, donnée qui dérange, coût de l'attente, ampleur du changement de taux ou désaccord avec le consensus. Appuyez chaque question sur un élément précis de la présentation ou une information vérifiable, pour amener l'équipe à préciser ou à défendre son raisonnement.

Préparez deux questions à partir des diapositives et adaptez-les aux explications entendues. \textbf{La seconde peut rebondir sur la première réponse; elle compte alors comme votre deuxième et dernière question.} Posez des questions courtes, sans accumulation de sous-questions. Il n'y a pas de question supplémentaire après la nouvelle. Aucun document supplémentaire n'est à remettre.

\textbf{Jouez le jeu!} Vous pouvez vous présenter au nom d'un \textbf{média réel ou fictif de votre choix}. Un nom inventif et une touche d'humour sont bienvenus. Cette brève introduction fait partie des quatre minutes de la conférence de presse.

Vos questions doivent être \textbf{exigeantes et respectueuses}. Leur qualité économique et votre écoute sont évaluées, indépendamment des réponses de l'autre équipe.

\newpage
\subsection{La nouvelle de dernière minute}
Après les deux questions des journalistes, le professeur présentera une \textbf{nouvelle économique hypothétique}, à traiter comme une information reçue avant la décision visée. Elle pourra concerner, par exemple, les prix, le marché du travail, l'activité économique ou la conjoncture internationale. Les équipes recevront des nouvelles différentes, de difficulté comparable.

Le scénario fournira le contexte et un point de comparaison pertinent pour interpréter la nouvelle. Après votre délibération, l'équipe présentatrice dispose de quatre minutes pour traiter les trois points suivants :
\begin{enumerate}
  \item \textbf{La décision prévue.} Annoncez si votre prévision change et justifiez son maintien ou sa révision.
  \item \textbf{Le degré de confiance.} Expliquez s'il change. Une nouvelle peut modifier les risques ou l'horizon d'un ajustement sans changer la décision que vous jugez la plus probable pour décembre.
  \item \textbf{Les conséquences pour votre industrie.} Distinguez les effets directs de la nouvelle de ceux d'une éventuelle réaction de la BdC. Expliquez si la nouvelle modifie votre analyse de la décision de gestion identifiée.
\end{enumerate}

\begin{pointcle}{Changer d'avis n'est pas une obligation}
Une réponse qui maintient ou renforce la prévision initiale peut obtenir tous les points. L'évaluation porte sur l'interprétation de la nouvelle, la cohérence du raisonnement et l'explication des incertitudes.
\end{pointcle}

\subsection{Prise de parole et utilisation de l'IA}
L'équipe répartit librement la prise de parole pendant l'exposé et les questions. Tous les membres doivent maîtriser l'analyse et contribuer à la préparation du rôle de journalistes. Le professeur modère les échanges et désignera \textbf{au moins un membre de l'équipe présentatrice pour répondre à une question}. Les réponses font partie de l'évaluation collective.

\textbf{La lecture d'un texte rédigé est interdite pendant l'exposé et les réponses}, sur papier ou à l'écran. Diapositives et mots-clés peuvent servir de repères, mais vous devez expliquer votre analyse dans vos propres mots. Les journalistes peuvent consulter leurs notes de préparation pour poser leurs questions.

L'IA est autorisée pour \textbf{préparer l'analyse, la présentation et les questions des journalistes}, mais \textbf{interdite pendant toute l'activité en classe, pour les deux rôles}, comme toute recherche externe ou aide extérieure. Vous demeurez responsables des données, calculs, références et arguments présentés.

\newpage
\section{Remise, évaluation et ressources}
\label{modalites}

\subsection{Diapositives à remettre}
Remettez \textbf{vos diapositives en format PDF au plus tard le 29 novembre 2026 à 23\,h\,59} (heure de Montréal). Elles doivent indiquer les membres de l'équipe, l'industrie, la variation exacte prévue du taux directeur, trois arguments clés et votre degré de confiance.

Identifiez la source et la période des données. Ajoutez à la fin des diapositives les références et la déclaration d'utilisation de l'IA : outils utilisés et leur rôle, ou absence d'utilisation. Distinguez les faits, les observations professionnelles et les hypothèses. Une prévision d'analyste peut servir de comparaison, mais ne remplace pas votre argument.

\subsection{Barème proposé}
La note comprend \textbf{90 points communs à l'équipe} et \textbf{10 points de contribution individuelle}.
\begingroup
\small
\begin{tabularx}{\textwidth}{@{}Y>{\raggedleft\arraybackslash}p{1.1cm}@{}}
\toprule
\textbf{Critère} & \textbf{Points} \\
\midrule
\textbf{Diagnostic et prévision} : qualité de l'analyse, interprétation et mécanismes économiques. & 20 \\
\addlinespace[4pt]
\textbf{Application à l'industrie} : utilisation pertinente du \textit{Rapport sur la politique monétaire}, justification de la sensibilité de l'industrie, qualité des mécanismes reliant le diagnostic aux effets attendus et à une décision de gestion précise. & 30 \\
\addlinespace[4pt]
\textbf{Réaction à la nouvelle} : interprétation de l'information, qualité du raisonnement économique, justification du maintien ou de la révision de la prévision et du degré de confiance, et implications pour l'industrie. & 20 \\
\addlinespace[4pt]
\textbf{Communication et réponses} : clarté de la présentation, force de l'argumentation, sources, qualité des réponses aux journalistes et respect du temps. & 10 \\
\addlinespace[4pt]
\textbf{Rôle de journalistes} : pertinence et précision des deux questions sur la prévision de décision de la BdC, ancrage dans l'analyse présentée et adaptation aux explications entendues. & 10 \\
\addlinespace[4pt]
\textbf{Contribution individuelle} : évaluations confidentielles par les pairs, examinées par le professeur. & 10 \\
\midrule
\textbf{Total avant bonus} & \textbf{100} \\
\bottomrule
\end{tabularx}
\endgroup

\subsection{Évaluation des contributions individuelles}
À l'issue des présentations, chaque membre remet au professeur une évaluation confidentielle de ses coéquipiers : contribution à l'analyse, aux diapositives et à la préparation du rôle de journalistes, fiabilité, respect des engagements et collaboration. Appuyez vos appréciations sur des exemples concrets et décrivez brièvement votre propre contribution. L'évaluation porte uniquement sur ce projet. Le professeur attribue les 10 points individuels à partir de ces appréciations et de leurs justifications; il ne s'agit pas d'une moyenne automatique des notes données par les pairs.

L'analyse est évaluée à partir des informations disponibles à la date limite de remise. \textbf{Une prévision qui ne se réalise pas peut obtenir tous les points du barème}; l'exactitude de la décision prévue intervient uniquement dans le bonus.

\textbf{Bonus de précision : 5 points sur la note du projet} si la variation exacte prévue correspond à celle annoncée le 9 décembre 2026. La prévision inscrite dans les diapositives est figée à la date limite commune de remise, le \textbf{29 novembre à 23\,h\,59}. La réponse hypothétique ne la modifie pas. Le bonus s'ajoute à la note de chaque membre, jusqu'à concurrence de 100 sur 100.

\end{document}
